% =====================================================================
% ЧАСТЬ III: Анализ траекторий и сценариев
% =====================================================================

\section{Три параллельных аналитика (Стадия 3)}

% Source: docs/04-analysts.md, Stage 3 specification

На Стадии 3 события, выделенные EventTrendAnalyzer, поступают в три специализированные аналитические агента, работающие \emph{параллельно}:

\begin{enumerate}
    \item \textbf{Геополитический аналитик} (\emph{Geopolitical Analyst}) — анализ стратегических акторов, баланса сил, альянсов, вероятности эскалации;
    \item \textbf{Экономический аналитик} (\emph{Economic Analyst}) — анализ финансовых потоков, рыночных индикаторов, цепочек поставок, сценариев воздействия на индексы;
    \item \textbf{Медиа-аналитик} (\emph{Media Analyst}) — оценка медийной ценности события через шести-мерный фреймворк \emph{Galtung \& Ruge} (1965).
\end{enumerate}

Минимум успешных оценок: $\text{min\_successful} = 2$ из 3. Архитектура допускает отказ одного аналитика без отмены Дельфи-прогноза.

\subsection{Геополитический анализ}

Для каждого события строится профиль ключевых стратегических акторов:
\begin{itemize}
    \item \textbf{Стратегические акторы}: государства, альянсы (НАТО, ШОС), международные организации (ООН, ЕС).
    \item \textbf{Баланс сил}: относительная мощь, направления сдвига в расстановке (кто усиливается, кто ослабевает).
    \item \textbf{Вероятность эскалации}: для конфликтов, дипломатических кризисов — числовая оценка (0–1).
    \item \textbf{Эффекты второго порядка}: цепь причин-следствий, как действие одного актора спровоцирует ответ другого.
\end{itemize}

Аналитическая рамка сочетает неореализм (структура системы определяет поведение акторов) с конструктивизмом (восприятие угроз зависит от идентичности и нарратива).

\subsection{Экономический анализ}

Для каждого события:
\begin{itemize}
    \item \textbf{Затронутые индикаторы}: валютные пары, цены на сырьё, фондовые индексы, облигационные спреды.
    \item \textbf{Фискальные последствия}: взаимосвязь с бюджетной политикой государств.
    \item \textbf{Цепочки поставок}: географические разрывы, сектора, уязвимые к эскалации или стабилизации.
    \item \textbf{Рыночный сигнал}: что рынки уже закладывают в цены и что ещё является сюрпризом.
\end{itemize}

Ключевой принцип: \emph{«Follow the money»} — экономические стимулы раскрывают истинные намерения и ограничения акторов.

\subsection{Медиа-анализ: шести-мерный фреймворк новостной ценности}

% Source: Galtung & Ruge 1965, Entman 1993 framing theory

Новостная ценность события оценивается по шести критериям (Galtung \& Ruge, 1965):

\begin{enumerate}
    \item \textbf{Timeliness} (своевременность): событие происходит сейчас, не архивное.
    \item \textbf{Magnitude} (масштаб): велико ли событие по количеству затронутых лиц или активов.
    \item \textbf{Prominence} (известность): вовлечены ли топ-фигуры, знакомые аудитории.
    \item \textbf{Proximity} (близость): географическая, культурная, тематическая близость к аудитории издания.
    \item \textbf{Conflict} (конфликтность): наличие явного противостояния, эмоционального заряда.
    \item \textbf{Novelty} (новизна): событие неожиданно, радикально отличается от прецедентов.
\end{enumerate}

Кроме того, медиа-аналитик оценивает:
\begin{itemize}
    \item \textbf{Редакционный fit} (0–1): соответствие события редакционной линии издания.
    \item \textbf{Медиа-насыщенность} (0–1): как долго тема находилась в топе новостей. Если > 14 дней подряд — редакция ищет свежий угол или отстранится.
    \item \textbf{Позиция в цикле}: находится ли тема на пике внимания, идёт на спад или восходящий тренд.
\end{itemize}

Выход: \texttt{MediaAssessment} с итоговым баллом новостной ценности и прогнозом вероятности попадания в выпуск данного издания.

---

% =====================================================================
% МОДЕЛИРОВАНИЕ ТРАЕКТОРИЙ СОБЫТИЙ
% =====================================================================

\section{Моделирование траекторий событий (Стадия 3)}

% Source: docs/04-analysts.md, EventTrajectory schema

После параллельного анализа для каждого события строится модель его траектории развития. Траектория описывает:

\begin{enumerate}
    \item \textbf{Текущее состояние} — где находится событие сейчас (описание 2–3 предложениями).
    \item \textbf{Моментум} (\emph{momentum}) — вектор развития события.
    \item \textbf{Три сценария} — BASELINE, OPTIMISTIC/PESSIMISTIC, WILDCARD.
    \item \textbf{Ключевые факторы} — 3–5 движущих сил, определяющих развитие.
    \item \textbf{Неопределённости} — 2–3 основные точки неопределённости.
\end{enumerate}

\subsection{Моментум события}

Моментум описывает вектор развития событийной нити в краткосрочной перспективе:

\begin{itemize}
    \item \textbf{Escalating} — ситуация накаляется, вероятность критических событий растёт.
    \item \textbf{Stable} — статус-кво поддерживается, существенных изменений не ожидается.
    \item \textbf{De-escalating} — напряжённость снижается, кризис миновал пик.
    \item \textbf{Emerging} — новое событие только начинает появляться в медийном пространстве.
    \item \textbf{Culminating} — событие приближается к критической точке, решению или кульминации.
    \item \textbf{Fading} — событие теряет актуальность, медийное внимание падает.
\end{itemize}

\subsection{Три сценария развития}

Для каждого события определяются три сценария с указанием:
\begin{itemize}
    \item Краткое описание (2–3 предложения).
    \item Назначенная вероятность (0.0–1.0), сумма по трём сценариям = 1.0.
    \item 2–3 ключевых индикатора, указывающих на реализацию этого сценария.
    \item Потенциальный заголовок, который может породить сценарий.
\end{itemize}

\textbf{Типы сценариев}:
\begin{enumerate}
    \item \textbf{BASELINE} — наиболее вероятное развитие событий на основе текущего моментума.
    \item \textbf{OPTIMISTIC} или \textbf{PESSIMISTIC} — улучшение или ухудшение ситуации по ключевым параметрам.
    \item \textbf{WILDCARD} — неожиданный поворот, требующий активации какого-либо из выявленных рисков.
\end{enumerate}

\subsection{Матрица перекрёстных влияний}

% Source: docs/04-analysts.md, CrossImpactMatrix schema

События в информационном пространстве не независимы. Развитие одного события влияет на вероятность других.

Матрица перекрёстных влияний строится как разреженное представление:
\begin{equation}
    \text{CrossImpactEntry:} \quad (source\_id, target\_id, impact\_score) \in [-1, 1] \times \mathbb{R}
\end{equation}

где:
\begin{itemize}
    \item $source\_id$ — событие-причина.
    \item $target\_id$ — событие-следствие.
    \item $impact\_score \in [-1, 1]$ — сила и направление влияния:
    \begin{itemize}
        \item $> 0$ — усиливающее влияние (источник ускоряет, вероятнее делает целевое событие).
        \item $< 0$ — ослабляющее влияние (источник замедляет, маловероятнее делает целевое).
        \item $= 0$ — отсутствие влияния (не включается в матрицу).
    \end{itemize}
\end{itemize}

Для типичного портфеля из 20 событий полная матрица содержала бы $20 \times 19 = 380$ пар. На практике значимых связей 30–50, остальные игнорируются. Эта разреженность критична для масштабируемости Дельфи при большом числе событий.

---

% =====================================================================
% ЧАСТЬ IV: МЕТОД ДЕЛЬФИ ДЛЯ LLM-АГЕНТОВ
% =====================================================================

\section{Пять экспертных персон (Дельфи)}

% Source: src/agents/forecasters/personas.py, lines 24-242

Дельфи-система построена на пяти независимых \emph{экспертных персонах} (\emph{expert personas}) — LLM-агентов с чётко определёнными аналитическими рамками, когнитивными смещениями и начальными весами. Каждая персона работает на одной и той же модели (\texttt{claude-opus-4.6}), но различается \emph{системным промптом}, определяющим её идентичность, аналитическую методологию и осознанные когнитивные смещения.

\begin{table}[h]
\small
\centering
\begin{tabular}{|c|l|l|c|l|}
\hline
\textbf{ID} & \textbf{Имя} & \textbf{Роль} & \textbf{Вес} & \textbf{Основная методология} \\
\hline
\hline
1 & REALIST & Аналитик рисков & 0.22 & Базовые ставки, Tetlock \emph{Superforecasting} \\
\hline
2 & GEOSTRATEG & Специалист МО & 0.20 & Неореализм, cui bono, баланс сил \\
\hline
3 & ECONOMIST & Макроэкономист & 0.20 & Follow the money, рыночные индикаторы \\
\hline
4 & MEDIA\_EXPERT & Редактор/медиа-аналитик & 0.18 & Гейткипинг \emph{White 1950}, фрейминг \emph{Entman 1993} \\
\hline
5 & DEVILS\_ADVOCATE & Контрариан/red teamer & 0.20 & Pre-mortem, \emph{Black swan} detection, Taleb \\
\hline
\end{tabular}
\caption{Экспертные персоны Дельфи: идентичность, методология и начальные веса}
\label{table:personas}
\end{table}

Суммарный вес: $0.22 + 0.20 + 0.20 + 0.18 + 0.20 = 1.00$.

\subsection{Персона 1: Реалист-аналитик}

\textbf{Идентичность}: Опытный аналитик политических рисков (условный профиль: Eurasia Group, Oxford Analytica) с 20-летним опытом работы в консалтинге. Мыслит категориями базовых ставок, исторических прецедентов и институциональной инерции.

\textbf{Аналитическая рамка}:
\begin{enumerate}
    \item \emph{Базовые ставки} — первый вопрос: «Как часто подобное происходило за 10–20 лет? Каковы были исходы?»
    \item \emph{Институциональная инерция} — системы меняются медленнее, чем кажется.
    \item \emph{Outside view перед inside view} — начинать с внешнего вида (Tetlock, Superforecasting).
    \item \emph{Калибровка над убеждённостью} — лучше неуверенная точная оценка, чем уверенная ошибочная.
    \item \emph{Разделение вопросов} — «Произойдёт ли?» и «Попадёт ли в заголовки?» — разные задачи.
\end{enumerate}

\textbf{Когнитивное смещение} (управляемое):
\begin{itemize}
    \item \textbf{Переоценивает}: инерция статус-кво, предсказуемость бюрократических процессов.
    \item \textbf{Недооценивает}: чёрные лебеди, скорость эскалации, роль личных решений лидеров.
    \item \textbf{Якорь при неопределённости}: исторические прецеденты.
\end{itemize}

\textbf{Начальный вес}: 0.22 (повышенный — базовые ставки исторически хорошо калиброваны).

\subsection{Персона 2: Геополитический стратег}

\textbf{Идентичность}: Специалист по международным отношениям (IISS, Chatham House, ИМЭМО РАН) с фокусом на силовые балансы, альянсы и стратегические интересы государств.

\textbf{Аналитическая рамка}:
\begin{enumerate}
    \item \emph{Cui bono} — кто выигрывает, кто проигрывает?
    \item \emph{Деревья решений} — эффекты второго и третьего порядка.
    \item \emph{Неореализм + конструктивизм} — структура системы определяет поведение; идентичности и нарративы формируют восприятие.
    \item \emph{Красные линии и пороги} — различение реальных красных линий от блефа.
\end{enumerate}

\textbf{Когнитивное смещение} (управляемое):
\begin{itemize}
    \item \textbf{Переоценивает}: рациональность государственных акторов, значимость геополитических факторов.
    \item \textbf{Недооценивает}: внутриполитические случайности, роль технологий, экономическую мотивацию негосударственных акторов.
    \item \textbf{Якорь}: модель великодержавного соперничества.
\end{itemize}

\textbf{Начальный вес}: 0.20.

\subsection{Персона 3: Экономический аналитик}

\textbf{Идентичность}: Макроэкономист и рыночный аналитик (Goldman Sachs Research, The Economist Intelligence Unit) с фокусом на потоки капитала, фискальную политику, цены на сырьё и корпоративные интересы.

\textbf{Аналитическая рамка}:
\begin{enumerate}
    \item \emph{Follow the money} — базовый принцип.
    \item \emph{Рациональный актор с бюджетными ограничениями} — экономические стимулы раскрывают истинные намерения.
    \item \emph{Экономический календарь} — источник предсказуемых новостных поводов (релизы, заседания ЦБ).
    \item \emph{Рыночные сигналы} содержат информацию (кривая доходностей, CDS спреды, цена золота).
\end{enumerate}

\textbf{Когнитивное смещение} (управляемое):
\begin{itemize}
    \item \textbf{Переоценивает}: экономическую рациональность, предсказательную силу рыночных индикаторов.
    \item \textbf{Недооценивает}: идеологические факторы, культурные войны, события, обусловленные эмоциями.
    \item \textbf{Якорь}: экономический календарь и рыночный консенсус.
\end{itemize}

\textbf{Начальный вес}: 0.20.

\subsection{Персона 4: Медиа-эксперт}

\textbf{Идентичность}: Бывший редактор крупного новостного агентства (Reuters, ТАСС) + преподаватель журналистики. Понимает внутреннюю кухню редакции: дедлайны, источники, редакционную политику, конкуренцию за внимание.

\textbf{Аналитическая рамка}:
\begin{enumerate}
    \item \emph{Гейткипинг} (\textbf{White}, 1950) — не всё важное публикуется. Редакционные фильтры ограничивают пропускную способность.
    \item \emph{Фрейминг} (\textbf{Entman}, 1993) — одно событие, десять подач. Выбор фрейма определяет заголовок и угол.
    \item \emph{Медиа-насыщенность} — тема в топе 14+ дней подряд $\to$ издание ищет свежий угол или отстранится.
    \item \emph{Конкуренция за внимание} — параллельные события конкурируют за аудиторию.
\end{enumerate}

\textbf{Когнитивное смещение} (управляемое):
\begin{itemize}
    \item \textbf{Переоценивает}: значимость медиа-циклов и информационных трендов, предсказуемость редакционных решений.
    \item \textbf{Недооценивает}: технические и процедурные события, медленно развивающиеся кризисы.
    \item \textbf{Якорь}: текущий новостной цикл и тематический баланс издания.
\end{itemize}

\textbf{Начальный вес}: 0.18 (ниже среднего — медиа-экспертиза важна для формулировки заголовков, менее для вероятностных оценок).

\subsection{Персона 5: Адвокат дьявола}

\textbf{Идентичность}: Систематический контрариан и специалист по анализу рисков (red team, разведывательное сообщество, школа Nassim Taleb). Задача — найти то, что остальные пропустили.

\textbf{Аналитическая рамка}:
\begin{enumerate}
    \item \emph{Pre-mortem} — «Представь, прогноз провалился. Что именно пошло не так?» (Klein, 1989).
    \item \emph{Steelmanning} — строить максимально сильную версию противоположной позиции, потом атаковать.
    \item \emph{Каскадные зависимости} — если A спорно и B зависит от A, то цепочка становится хрупкой.
    \item \emph{Чёрные лебеди} (Талеб) — маловероятные, но высокоимпактные сценарии.
\end{enumerate}

\textbf{Когнитивное смещение} (управляемое, и в этом случае — желаемое):
\begin{itemize}
    \item \textbf{Переоценивает}: вероятность чёрных лебедей, хрупкость систем, неочевидные причинно-следственные связи.
    \item \textbf{Недооценивает}: устойчивость статус-кво, вероятность «скучных» исходов, стабильность институтов.
    \item \textbf{Якорь}: маловероятные, но высокоимпактные сценарии.
\end{itemize}

\textbf{Начальный вес}: 0.20 (парадоксально высокий — контрарианские прогнозы добавляют ценность ансамблю при агрегации, даже если редко верны в большинстве случаев).

---

% =====================================================================
% РАУНД 1: НЕЗАВИСИМЫЕ ОЦЕНКИ
% =====================================================================

\section{Раунд 1: независимые оценки пяти персон}

% Source: src/agents/forecasters/personas.py, lines 269-333, and docs/05-delphi-pipeline.md §2

Все пять персон запускаются \emph{параллельно}. Каждая получает:

\begin{itemize}
    \item \textbf{Профиль издания} (\texttt{OutletProfile}): географический охват, целевая аудитория, редакционная политика, тематический баланс.
    \item \textbf{Траектории событий} (\texttt{EventTrajectory[]}): до 20 событий с текущим состоянием, моментумом, сценариями, ключевыми факторами.
    \item \textbf{Матрица перекрёстных влияний} (\texttt{CrossImpactMatrix}): разреженное представление зависимостей между событиями.
    \item \textbf{Система промпта} — специфичный для персоны системный промпт, определяющий её идентичность, методологию и смещения.
\end{itemize}

Параметры LLM для R1:
\begin{itemize}
    \item \textbf{Модель}: \texttt{anthropic/claude-opus-4.6}
    \item \textbf{Temperature}: $T = 0.7$ (некоторая стохастичность, но управляемая)
    \item \textbf{max\_tokens}: 4096
    \item \textbf{json\_mode}: истина (структурированный выход)
\end{itemize}

Минимум успешных агентов: $\text{min\_successful} = 3$ из 5. При отказе одного агента раунд продолжается.

\subsection{Выходная схема: PersonaAssessment}

Каждая персона возвращает структурированную оценку \texttt{PersonaAssessment}:

\begin{definition}[\texttt{PersonaAssessment} — оценка от одной персоны]
\begin{verbatim}
{
  persona_id: str,               # "realist", "geostrateg", и т.д.
  round_number: int,             # 1 или 2
  predictions: PredictionItem[], # 5–15 прогнозов
  cross_impacts_noted: str[],    # замеченные перекрёстные влияния
  blind_spots: str[],            # что группа может пропустить
  confidence_self_assessment: float,  # (0–1) мета-уверенность

  # Только для раунда 2:
  revisions_made: str[],         # что было изменено
  revision_rationale: str        # почему
}
\end{verbatim}
\end{definition}

\subsection{Выходная схема: PredictionItem}

Каждый элемент \texttt{predictions[]} содержит одно предсказание события:

\begin{definition}[\texttt{PredictionItem} — единичный прогноз]
\begin{verbatim}
{
  event_thread_id: str,              # ID события из Stage 3
  prediction: str,                   # конкретное утверждение
  probability: float,                # вероятность (0.0–1.0, не кратна 5%)
  newsworthiness: float,             # (0–1) вероятность попадания в выпуск
  scenario_type: ScenarioType,       # BASELINE | OPTIMISTIC | PESSIMISTIC | WILDCARD
  reasoning: str,                    # цепочка рассуждений (3–7 предложений)
  key_assumptions: str[],            # 2–4 ключевые предпосылки
  evidence: str[],                   # ссылки на факты из входных данных
  conditional_on: str[]              # ID других прогнозов, от которых зависит этот
}
\end{verbatim}
\end{definition}

\textbf{Критические требования}:
\begin{enumerate}
    \item \textbf{Точность вероятностей}: каждая вероятность должна быть уникальна, не кратна 5 или 10 (0.63, не 0.60). Это требование (Tetlock) улучшает калибровку.
    \item \textbf{Диапазон}: $0.03 \leq \text{probability} \leq 0.97$. Вероятности 0.00 или 1.00 запрещены (они означают избыточную уверенность).
    \item \textbf{Обоснование}: каждый прогноз должен иметь явную цепочку рассуждений и ключевые предпосылки, на которых основана оценка.
\end{enumerate}

---

% =====================================================================
% МЕДИАТОР: СИНТЕЗ РАЗНОГЛАСИЙ
% =====================================================================

\section{Медиатор: синтез разногласий (Стадия 5a)}

% Source: src/agents/forecasters/mediator.py, lines 1-269; docs/prompts/mediator.md §1-10

После Раунда 1 пять независимых оценок поступают в \emph{Медиатор} — специализированный агент, чья задача \emph{не прогнозировать}, а \emph{структурировать разногласия} для содержательного второго раунда.

\subsection{Классический Дельфи vs. LLM-Дельфи}

В классическом методе Дельфи (Dalkey \& Helmer, 1963) обратная связь для второго раунда состоит из агрегированной статистики: медиана, квартили, гистограмма.

Однако исследование \textbf{DeLLMphi} (Zhao et al., 2024) показало: если LLM-агентам дать только медианные оценки группы, они слегка сдвигаются к ней без содержательной ревизии аргументов. Это явление называется \emph{Degeneration-of-Thought} (Liang et al., 2024, EMNLP).

\textbf{Наше решение}: вместо голой статистики, Медиатор формулирует \emph{содержательные конкретные вопросы}, которые раскрывают фактические разногласия и заставляют агентов пересматривать не числа, а аргументы.

\subsection{Три этапа медиации}

\subsubsection{Этап 1: Алгоритмическая классификация событий}

Для каждого события Медиатор собирает все оценки пяти персон и вычисляет:

\begin{enumerate}
    \item \textbf{Разброс} (\emph{spread}): $\text{spread} = \max(\text{probabilities}) - \min(\text{probabilities})$
    \item \textbf{Медиана}: $\text{median}(\text{probabilities})$
    \item \textbf{Число агентов, упомянувших событие}: $n \in [0, 5]$
\end{enumerate}

События классифицируются как:

\begin{enumerate}
    \item \textbf{Консенсус} (\emph{Consensus Area}): $\text{spread} < 0.15$ \textbf{AND} $n \geq 3$

    \textit{Интерпретация}: эксперты согласны, нет необходимости в ревизии. Это событие попадёт в финальный прогноз с доверием.

    \item \textbf{Разногласие} (\emph{Dispute Area}): $\text{spread} \geq 0.15$

    \textit{Интерпретация}: эксперты не согласны. Медиатор формулирует конкретный фактический вопрос, ответ на который мог бы сблизить позиции.

    \item \textbf{Пробел} (\emph{Gap Area}): $n < 3$

    \textit{Интерпретация}: события упомянуты менее чем тремя экспертами. Потенциально важные события, которые группа в целом упустила.
\end{enumerate}

\subsubsection{Этап 2: Проверка каскадных зависимостей}

\emph{CrossImpactFlag}: если прогноз события A зависит от события B (через поле \texttt{conditional\_on}), и по B существует разногласие, то это создаёт цепную неопределённость, которую нужно подсветить.

\subsubsection{Этап 3: LLM-обогащение синтеза}

Медиатор получает:
\begin{itemize}
    \item Анонимизированные оценки (метки Эксперт A–E, см. ниже).
    \item Траектории событий (для контекста).
    \item Результаты этапов 1–2 (алгоритмическая классификация).
\end{itemize}

И через LLM-вызов \emph{обогащает} синтез:
\begin{itemize}
    \item Для каждого разногласия: формулирует one-liner ключевого вопроса, ответ на который сближает позиции.
    \item Для каждого пробела: объясняет, почему это может быть важным упущением.
    \item Генерирует общее резюме (2–3 предложения): сколько консенсусов, расхождений, пробелов.
\end{itemize}

Параметры LLM для Медиатора:
\begin{itemize}
    \item \textbf{Модель}: \texttt{anthropic/claude-opus-4.6}
    \item \textbf{Temperature}: 0.7 (нейтральная)
    \item \textbf{json\_mode}: истина
\end{itemize}

\subsection{Анонимность как защита от группового мышления}

% Source: src/agents/forecasters/mediator.py, lines 181-206; docs/prompts/mediator.md §7-9

Критический механизм защиты от давления большинства (Zhang et al., 2024, ACL): **метки экспертов переиграются случайно при каждом запуске**.

Алгоритм:
\begin{enumerate}
    \item Для каждого набора оценок генерируется детерминированное дерево случайных чисел, основанное на хеше идентификаторов персон.
    \item Метки $\{$Эксперт A, B, C, D, E$\}$ перемешиваются.
    \item Одна и та же персона получает разную метку при каждом прогоне.
\end{enumerate}

Эффект: агент не может узнать себя по метке и подстраиться под неё. Это предотвращает конформность и сохраняет независимость позиций меньшинства.

\subsection{Выходная схема: MediatorSynthesis}

\begin{definition}[\texttt{MediatorSynthesis} — структурированный синтез разногласий]
\begin{verbatim}
{
  consensus_areas: ConsensusArea[],     # spread < 0.15, n >= 3
  disputes: DisputeArea[],              # spread >= 0.15 + ключевой вопрос
  gaps: GapArea[],                      # n < 3
  cross_impact_flags: CrossImpactFlag[], # зависимости между спорными событиями
  overall_summary: str,                 # 2–3 предложения об итоговой картине
  supplementary_facts: str[]            # заполняется позже, если нужен supervisor_search
}
\end{verbatim}
\end{definition}

где:

\begin{definition}[\texttt{ConsensusArea}]
\begin{verbatim}
{
  event_thread_id: str,
  median_probability: float,  # (0–1)
  spread: float,              # < 0.15
  num_agents: int             # >= 3
}
\end{verbatim}
\end{definition}

\begin{definition}[\texttt{DisputeArea}]
\begin{verbatim}
{
  event_thread_id: str,
  median_probability: float,
  spread: float,              # >= 0.15
  positions: AnonymizedPosition[],  # позиции каждого эксперта
  key_question: str           # конкретный фактический вопрос для R2
}

AnonymizedPosition: {
  agent_label: str,           # "Эксперт A", "Эксперт B", и т.д.
  probability: float,
  reasoning_summary: str,     # первые 200 символов обоснования
  key_assumptions: str[]
}
\end{verbatim}
\end{definition}

\begin{definition}[\texttt{GapArea}]
\begin{verbatim}
{
  event_thread_id: str,
  mentioned_by: str[],        # список меток экспертов, упомянувших событие
  note: str                   # почему это может быть важным пробелом
}
\end{verbatim}
\end{definition}

\begin{definition}[\texttt{CrossImpactFlag}]
\begin{verbatim}
{
  prediction_event_id: str,    # событие A
  depends_on_event_id: str,    # событие B (спорное)
  note: str                    # описание зависимости A$\to$B
}
\end{verbatim}
\end{definition}

---

% =====================================================================
% РАУНД 2: РЕВИЗИЯ С ОБРАТНОЙ СВЯЗЬЮ
% =====================================================================

\section{Раунд 2: ревизия с обратной связью}

% Source: src/agents/forecasters/personas.py, lines 269-308; docs/05-delphi-pipeline.md §3.2-3.3

Все пять персон запускаются \emph{снова параллельно}. Каждая получает:

\begin{itemize}
    \item \textbf{Свои собственные оценки из R1}.
    \item \textbf{MediatorSynthesis} — анонимизированные позиции других экспертов и ключевые вопросы для разногласий.
    \item \textbf{Гвард независимости} (\emph{Independence Guard}): явная инструкция — \textit{«НЕ сдвигай числа просто потому, что все думают иначе. Ответь на ключевой вопрос медиатора аргументированно»}.
\end{itemize}

Параметры LLM для R2:
\begin{itemize}
    \item \textbf{Модель}: \texttt{anthropic/claude-opus-4.6}
    \item \textbf{Temperature}: $T = 0.6$ (снижена с 0.7, чтобы уменьшить шум, но сохранить творчество)
    \item \textbf{max\_tokens}: 4096
    \item \textbf{json\_mode}: истина
\end{itemize}

Минимум успешных агентов: $\text{min\_successful} = 3$ из 5.

\subsection{Выходная схема: RevisedAssessment}

Структура идентична \texttt{PersonaAssessment}, но с дополнительными полями:

\begin{verbatim}
{
  # Все поля как в R1...
  predictions: PredictionItem[],
  confidence_self_assessment: float,

  # Плюс, для R2:
  revisions_made: str[],        # список ревизий
  revision_rationale: str       # почему эти ревизии
}
\end{verbatim}

Примеры ревизий:
\begin{itemize}
    \item «Повысил вероятность с 0.42 до 0.58 — ключевой вопрос медиатора о венгерском вето важен, я пересмотрел рассуждения».
    \item «Сохранил 0.25 — позиция меньшинства, но фактический вопрос не рассеял мою неопределённость».
    \item «Добавил новый прогноз на событие, упомянутое одним экспертом — теперь вижу важность этого пробела».
\end{itemize}

---

% =====================================================================
% ТЕОРЕТИЧЕСКОЕ ОБОСНОВАНИЕ МЕТОДА ДЕЛЬФИ
% =====================================================================

\section{Теоретическое обоснование метода Дельфи для LLM-агентов}

\subsection{Классический Дельфи и его преимущества}

\emph{Метод Дельфи} (Dalkey \& Helmer, RAND Corporation, 1963) — структурированная техника группового прогнозирования, основанная на четырёх принципах:

\begin{enumerate}
    \item \textbf{Анонимность}: эксперты не знают, кто дал какую оценку.
    \item \textbf{Итерация}: несколько раундов оценок с обратной связью.
    \item \textbf{Контролируемая обратная связь}: участники получают агрегированную статистику группы (медиана, квартили, гистограмма).
    \item \textbf{Статистическая агрегация}: финальный результат — усреднённая оценка группы.
\end{enumerate}

Доказанные преимущества (Rowe \& Wright, 2001, 2005):
\begin{itemize}
    \item Подавление эффекта якорения — эксперты сдвигают позиции на основе данных, а не авторитета.
    \item Снижение давления авторитета — анонимность защищает меньшинство.
    \item Более калиброванные вероятности — групповая оценка часто превосходит одиночных экспертов.
\end{itemize}

\subsection{Адаптация для LLM-агентов}

Наша реализация заменяет человеческих экспертов на LLM-агентов с чётко определёнными когнитивными профилями. Ключевые отличия:

\begin{enumerate}
    \item \textbf{Детерминизм} — каждый агент получает одинаковые входные данные (траектории, перекрёстные влияния) и системный промпт, определяющий его аналитическую рамку.

    \item \textbf{Отсутствие давления авторитета в классическом смысле} — агенты не знают о мнениях друг друга до медиации. Давление группы вводится контролируемо через медиатора, а не через выраженное большинство.

    \item \textbf{Воспроизводимость} — результаты должны быть воспроизводимы (с фиксированными семенами случайности).
\end{enumerate}

\subsection{Ключевые исследования}

\subsubsection{AIA Forecaster: необходимость модельного разнообразия}

Schoenegger et al. (2024) показали: если все агенты в ансамбле используют одну модель, их ошибки \emph{коррелируют}. Коррелированные ошибки не компенсируются при агрегации.

\textbf{Вывод}: ансамбль из 5 экземпляров одной модели лишь незначительно лучше одного вызова той же модели.

\textbf{Наше решение}: все пять персон работают на одной модели (\texttt{claude-opus-4.6}), но \emph{разнообразие ошибок обеспечивается система prompts, когнитивными смещениями и аналитическими рамками}, не различием моделей. Это позволяет:
\begin{itemize}
    \item Контролировать стоимость (все вызовы к одному провайдеру).
    \item Гарантировать минимальное качество рассуждений (Opus 4.6 — лучший в своём ценовом диапазоне).
    \item Достичь разнообразия через когнитивные профили (исследовано).
\end{itemize}

\subsubsection{DeLLMphi и проблема Degeneration-of-Thought}

Zhao et al. (2024) воспроизвели классический Дelphi с LLM и обнаружили критическую проблему:

\begin{enumerate}
    \item В R1 агенты дают независимые оценки.
    \item В R2 им показывают медианные оценки группы.
    \item Агенты \emph{слегка сдвигаются к медиане}, но \textbf{без содержательной ревизии аргументов}.
\end{enumerate}

Это явление названо \emph{Degeneration-of-Thought} (Liang et al., 2024, EMNLP): LLM, утвердившиеся в позиции, не способны к \emph{genuine revision} через self-reflection одной только статистики.

\textbf{Наше решение}: Медиатор не просто даёт медиану. Он формулирует \emph{конкретные фактические вопросы}, на которые необходимо ответить:

\begin{equation}
    \text{Ключевой вопрос:} \quad Q = \text{«Что именно нужно проверить для разрешения этого спора?»}
\end{equation}

Lorenz \& Fritz (2025, arXiv:2602.08889) показали: когда агентам задаются \emph{substantive questions}, корреляция их оценок с ground truth достигает $r = 0.87$--$0.95$. При простой статистике $r \approx 0.60$--$0.70$.

\subsubsection{Защита позиций меньшинства}

Zhang et al. (2024, ACL) показали: LLM-агенты воспроизводят human-like conformity с \emph{bandwagon score} 0.524 в GPT-3.5. То есть, когда большинство согласно, агент сдвигается к большинству даже если изначально думал иначе.

\textbf{Наше решение}: анонимность с ротацией меток. Метки (A–E) переиграются при каждом прогоне так, что агент не может узнать себя и подстраиться.

\subsubsection{Два раунда достаточно}

Классические исследования Дельфи (Rowe \& Wright, 2005) показали:
\begin{itemize}
    \item Раунд 1 $\to$ Раунд 2: существенное улучшение (до 80% сходимости).
    \item Раунд 2 $\to$ Раунд 3: маргинальное улучшение (дополнительно ~5–10%).
    \item Раунд 3+: убывающая отдача.
\end{itemize}

При LLM стоимость растёт линейно с числом раундов. \textbf{Наша архитектура: ровно 2 раунда} — баланс между качеством и стоимостью.

---

% =====================================================================
% Заключение логической архитектуры Дельфи
% =====================================================================

\section{Логическая архитектура двухраундового Дельфи}

Полный цикл:

\begin{equation}
\begin{aligned}
    \text{R1:} \quad &\text{5 personas} \xrightarrow{\parallel} \text{PersonaAssessment}[] \\
    \text{Mediation:} \quad &\text{PersonaAssessment}[] \xrightarrow{\text{Mediator}} \text{MediatorSynthesis} \\
    \text{R2:} \quad &\{\text{Self R1} + \text{MediatorSynthesis}\} \xrightarrow{\parallel} \text{RevisedAssessment}[] \\
    \text{Judge:} \quad &\text{RevisedAssessment}[] \xrightarrow{\text{aggregation}} \text{RankedPrediction}[]
\end{aligned}
\end{equation}

Каждая стадия:
\begin{itemize}
    \item Независима по данным (агенты не видят выводов друг друга до медиации).
    \item Парсируется в структурированный Pydantic-объект для валидации.
    \item Логируется с полной стоимостью LLM (tokens, USD).
    \item Переносит цепочку рассуждений, а не только вероятности.
\end{itemize}

Полная спецификация Delphi: `docs/05-delphi-pipeline.md`.

